%                                                                 aa.dem
% AA vers. 9.1, LaTeX class for Astronomy & Astrophysics
% demonstration file
%                                                       (c) EDP Sciences
%-----------------------------------------------------------------------
%
%\documentclass[referee]{aa} % for a referee version
%\documentclass[onecolumn]{aa} % for a paper on 1 column  
%\documentclass[longauth]{aa} % for the long lists of affiliations 
%\documentclass[letter]{aa} % for the letters 
%\documentclass[bibyear]{aa} % if the references are not structured 
%                              according to the author-year natbib style

%
\documentclass{aa}  

\usepackage[utf8]{inputenc}
\usepackage{natbib,twoopt}
\usepackage[breaklinks=true]{hyperref} %% to avoid \citeads line fills
\bibpunct{(}{)}{;}{a}{}{,}             %% natbib format for A&A and ApJ
\usepackage{graphicx}
\usepackage{array}
\usepackage{txfonts}
\usepackage{hyperref}
\hypersetup{
    colorlinks=true,
    linkcolor=blue,
    citecolor = blue,
    filecolor=magenta,      
    urlcolor=blue}
%
% ---------- macros ----------
\newcommand{\angstrom}{\mbox{\normalfont\AA}}
\newcommand{\todo}[1]{\textcolor{red}{[#1]}}
\newcommand{\salvador}{salvador}   % name of the code: change here only


\authorrunning{Alvarez Garay et al.}
\titlerunning{---}

\begin{document}

%---------------------------------------------------------------------
\title{salvador}
\subtitle{A code for abundance determination through spectral synthesis}

\author{D.~A. Alvarez Garay\inst{1}\fnmsep\thanks{Corresponding author:
        \email{deimer.alvarez@inaf.it}}
   \and \todo{friends}
       }

\institute{Osservatorio Astrofisico di Arcetri, Largo E. Fermi 5,
           50125 Firenze, Italy
          }

%\date{Received ...; accepted ...}

%---------------------------------------------------------------------
\abstract
{BLABLA}
\keywords{techniques: spectroscopic --
          methods: data analysis --
          stars: abundances --
          stars: atmospheres --
          stars: rotation
          }

\maketitle

%=====================================================================
\section{Introduction}
\label{sec:intro}
%=====================================================================

The chemical composition of stellar photospheres is one of the most
powerful diagnostics in stellar and Galactic astrophysics. Chemical
abundances trace the nucleosynthetic history of the gas from which
stars formed, the mixing processes that modify the surface
composition along stellar evolution, and the self-enrichment
processes responsible for the multiple populations observed in
globular clusters \citep[GCs; e.g.,][]{Bastian2018, Gratton2019}. In
the last two decades, large spectroscopic surveys such as APOGEE
\citep{Majewski2017}, the Gaia-ESO Survey \citep{Gilmore2022}, and
GALAH \citep{Buder2021}, together with multi-object and integral-field
spectrographs such as FLAMES \citep{Pasquini2002} and MUSE
\citep{Bacon2010}, and very high-resolution spectrographs such as
ESPRESSO \citep{Pepe2021}, have provided spectra for hundreds of
thousands of stars, with resolving powers ranging from
$R\sim3000$ to more than 100\,000. This wealth of data requires
analysis tools that are automatic and reproducible, while being
flexible enough to deal with spectral features of very different
nature: weak and strong atomic lines, blended features, molecular
bands, and transitions affected by hyperfine or isotopic splitting.

Chemical abundances are commonly derived following two approaches.
In the first one, the equivalent widths (EWs) of isolated lines are
measured, either by direct integration or by fitting analytical
profiles \citep[e.g., with DAOSPEC or ARES;][]{Stetson2008,
Sousa2007, Sousa2015}, and compared with the theoretical curve of
growth computed for a given model atmosphere, as done for instance by
MOOG \citep{Sneden1973}, GALA \citep{Mucciarelli2013}, and FAMA
\citep{Magrini2013}. This approach is fast and naturally suited to the
spectroscopic derivation of the atmospheric parameters through the
excitation and ionization balance. However, it is restricted to
unblended lines with a well-defined profile, and its reliability
decreases at low and intermediate spectral resolution, in crowded
spectral regions (as in metal-rich or cool stars), for strong lines
with extended damping wings, and for features affected by hyperfine
or isotopic splitting or by molecular blends. In the second approach,
the observed spectrum is directly compared with synthetic spectra.
Several codes adopt this strategy, either by computing the synthetic
spectra on the fly, as in SME \citep{Valenti1996, Piskunov2017} and
its Python implementation PySME \citep{Wehrhahn2023}, the
synthesis driver of MOOG, BACCHUS \citep{Masseron2016} and TSFitPy
\citep{Gerber2023}, both based on Turbospectrum
\citep{AlvarezPlez1998, Plez2012}, iSpec \citep{BlancoCuaresma2014,
BlancoCuaresma2019}, and FASMA \citep{Tsantaki2018}, or by
interpolating within precomputed grids or libraries, as in MyGIsFOS
\citep{Sbordone2014}, SP\_Ace \citep{Boeche2016}, and the APOGEE
pipeline ASPCAP \citep{GarciaPerez2016}, which relies on the FERRE
code \citep{AllendePrieto2006}. Data-driven methods, such as The
Cannon \citep{Ness2015} and The Payne \citep{Ting2019}, represent a
further step towards the fast analysis of very large samples. Many of
these tools determine atmospheric parameters and chemical abundances
simultaneously, often by fitting extended spectral regions.

Despite this variety of tools, chemical abundances derived from the
same spectra by different codes and groups can differ significantly,
because of differences in the adopted model atmospheres, atomic data,
continuum placement, line selection, and treatment of the line
broadening \citep{Smiljanic2014, Hinkel2016, Jofre2019}. A full control
of each step of the analysis, line by line, therefore remains
essential, in particular when small abundance variations must be
detected, as in the case of the multiple populations of GCs, or when
the abundances rely on a few, possibly blended or saturated,
spectral features.

In this paper we present \salvador{}\footnote{The code is publicly
available at \url{https://github.com/DeimerAlvarez/salvador}.}, a
Python code designed to derive chemical abundances through spectral
synthesis on a line-by-line basis. For each line, \salvador{} performs
a $\chi^2$ minimization between the observed spectrum and a set of
synthetic spectra in which only the abundance of the element under
analysis is varied. The synthetic spectra are computed on the fly with
the SYNTHE code \citep{Kurucz1993b, Sbordone2004, Kurucz2005}, adopting
one-dimensional, plane-parallel, local thermodynamic equilibrium (LTE)
ATLAS9 model atmospheres \citep{Kurucz1993, Kurucz2005} computed with the
opacity distribution functions of the KOALA database
\citep{Mucciarelli2026}, both
managed through the Python code K2 (Mucciarelli et al., in prep.). Alternatively, \salvador{} can use a grid of
precomputed synthetic spectra, calculated with any spectral synthesis
code. \salvador{} does not derive the atmospheric parameters, which
must be known beforehand, for instance from photometry or from the
spectroscopic analysis described in Sect.~\ref{sec:gala}. The code is
designed for the analysis of FGK stars, both in the optical and in the
infrared domain, and it can handle spectra with resolving powers from
$R\sim3000$, as those provided by MUSE, up to the very high resolution
of ESPRESSO. Previous versions of \salvador{} have already been
employed to derive chemical abundances in the GCs $\omega$~Centauri
and M\,54 \citep{AlvarezGaray2024, AlvarezGaray2026}.

Besides the standard derivation of abundances from the line profiles,
\salvador{} implements a set of features that, to our knowledge, are
not combined within a single code:
\begin{enumerate}
\item the abundance variation can be mimicked by an equivalent
variation of the oscillator strength ($\log gf$) of the transitions of
the species under analysis. In this way, the line opacity is modified
while the chemical composition adopted in the synthesis, and therefore
the ionization and molecular equilibria, the electron pressure, and
the continuous opacity, remain those of the model atmosphere. The
synthetic spectrum is thus fully consistent with the adopted model
atmosphere, also for elements that are important electron donors
(e.g., Mg, Si, and Fe) or that take part in the molecular equilibrium
(e.g., C, N, and O);
\item the spectral resolution can be estimated automatically for each
line, by comparing the full width at half maximum of the observed
line with that of the intrinsic (unbroadened) synthetic profile. This
is particularly useful for multi-object and integral-field
spectrographs, where the resolving power can vary with wavelength and
across the detector;
\item the abundances of strong, saturated lines can be derived from
their wings only, excluding the line core, which is almost
insensitive to the abundance and can be affected by departures from
LTE;
\item isotopic ratios, such as $^{12}$C/$^{13}$C and
$^{6}$Li/$^{7}$Li, can be measured;
\item the projected rotational velocity ($v\sin i$) can be derived
from individual lines, either simultaneously with the abundance or at
a fixed abundance;
\item the uncertainties arising from the atmospheric parameters are
computed line by line by repeating the full analysis with the
perturbed parameters, and they are propagated keeping their sign, so
that the sensitivities common to the element under analysis and to
iron partially cancel out in the [X/Fe] abundance ratios. The internal
uncertainties are estimated through Monte Carlo simulations based on
the signal-to-noise ratio measured around each line.
\end{enumerate}
\salvador{} can be run both in an interactive mode, which allows the
user to inspect and, when needed, adjust every fit (normalization,
wavelength shift, fitting region, and resolution), and in a fully
automatic mode, in which several lines can be analyzed in parallel.

Although \salvador{} is designed and tested for FGK stars, it can in
principle be applied also to M giants, since the TiO and H$_2$O line
lists, which dominate the opacity in the spectra of the coolest
stars, can be included in the computation of the synthetic spectra.
However, because of the huge number of molecular transitions involved,
the inclusion of these line lists makes the computation of each
synthetic spectrum very time consuming, and we therefore recommend
this option only for specific cases (Sect.~\ref{sec:synth_k2}).

This paper is organized as follows. In Sect.~\ref{sec:overview} we
give an overview of \salvador{} and of its input files and running
modes. In Sect.~\ref{sec:synth} we describe the computation of the
synthetic spectra. In Sect.~\ref{sec:method} we detail the analysis of
the individual lines, and in Sect.~\ref{sec:errors} the estimate of
the uncertainties. In Sect.~\ref{sec:implementation} we describe the
implementation and the performance of the code. In
Sect.~\ref{sec:gala} we present the new version of GALA, which
derives the atmospheric parameters. In Sect.~\ref{sec:validation} we
validate \salvador{} against synthetic spectra and benchmark stars,
and in Sect.~\ref{sec:conclusions} we summarize our work.

%=====================================================================
\section{Overview of the code}
\label{sec:overview}
%=====================================================================

\salvador{} is designed to measure, for each star, a list of spectral
features defined by the user, keeping the analysis of every feature
independent of the others. The main steps of the analysis are
summarized in Fig.~\ref{fig:flowchart}. For each observed spectrum,
the user provides the atmospheric parameters (effective temperature
$T_{\rm eff}$, surface gravity $\log g$, metallicity [M/H], and
microturbulent velocity $\xi$), the chemical mixture of the model
atmosphere, the radial velocity (RV) of the star, and a line list
specifying, for each feature, the kind of analysis to be performed,
the wavelength and the species of the transition, the spectral
resolution, the limits of the region where the fit is performed, and
the initial abundance. The chemical mixture is chosen among five
values of the $\alpha$-element enhancement, $[\alpha/{\rm
Fe}]=-0.4$, $-0.2$, $0.0$, $+0.2$, and $+0.4$\,dex, matching the
sampling of the KOALA grid of ATLAS9 opacity distribution functions
\citep[ODFs;][]{Mucciarelli2026}. Optionally, the uncertainties of
$T_{\rm eff}$, $\log g$, and $\xi$ can be given, in which case the
corresponding uncertainties on the abundances are computed
(Sect.~\ref{sec:errors_params}).

The kind of analysis is set for each feature by a keyword, listed in
Table~\ref{tab:keywords}. All the keywords share the same core
(Sect.~\ref{sec:method}), that is, the comparison between the observed
feature and a set of synthetic spectra that differ only in the quantity
to be measured: the abundance of the element (or, equivalently, the
$\log gf$ of its transitions), an isotopic ratio, or the projected
rotational velocity. A file of abundance variations, associated with
each spectrum, allows the user to change the abundance of any element
with respect to the chemical mixture of the model atmosphere, or the
isotopic composition of a species, in all the synthetic spectra of
that star. This is useful, for instance, to include in the synthesis
the known abundances of the elements that contribute to the blends of
a given feature, as done by \citet{AlvarezGaray2024}, who fixed the N
abundance when measuring Al from the doublet at 7835--7836\,\AA,
contaminated by CN lines.

\salvador{} can be run in two modes. In the automatic mode, all the
features are analyzed without any intervention by the user, and
several features can be analyzed in parallel
(Sect.~\ref{sec:implementation}). In the interactive mode, the fit of
each feature is displayed and the user can accept or reject it, or
modify the ingredients of the fit (Sect.~\ref{sec:interactive}). In
both modes, the synthetic spectra can be computed on the fly
(Sect.~\ref{sec:synth_k2}) or taken from a grid of precomputed
synthetic spectra (Sect.~\ref{sec:synth_grid}).

For each spectrum, \salvador{} provides a file with the results of
each feature (abundance, EW, excitation potential $\chi$, $\log gf$,
and the uncertainties described in Sect.~\ref{sec:errors}), a
document with the plots of all the fits, and diagnostic plots of the
Fe abundances as a function of $\chi$, reduced EW, and wavelength,
useful to check the adopted $T_{\rm eff}$ and $\xi$. The final
abundance ratios, [Fe/H] and [X/Fe], of all the stars and their
uncertainties are collected in a single table. The solar reference
abundances can be chosen among those by \citet{Caffau2011} and
\citet{Lodders2010}, which is the default and the reference mixture
of the KOALA models \citep{Mucciarelli2026}, \citet{Grevesse1998},
\citet{Asplund2009}, and \citet{Magg2022}.

\begin{figure}
\centering
\fbox{\parbox[c][5cm][c]{0.9\hsize}{\centering
\todo{Fig. 1: flowchart of salvador (input files $\rightarrow$
normalization and local RV $\rightarrow$ synthetic spectra $\rightarrow$
$\chi^2$ minimization $\rightarrow$ EW, Monte Carlo, parameter
reruns $\rightarrow$ final abundances)}}}
\caption{Schematic workflow of \salvador.}
\label{fig:flowchart}
\end{figure}

\begin{table}
\caption{Keywords defining the kind of analysis of each feature.}
\label{tab:keywords}
\centering
\begin{tabular}{l >{\raggedright\arraybackslash}p{6.0cm}}
\hline\hline
Keyword & Analysis \\
\hline
ABU    & Abundance from the line profile (Sect.~\ref{sec:abu}). \\
ABUF   & As ABU, but the EW includes all the transitions of the species
         in the feature (e.g., the Li doublet at 6708\,\AA). \\
ABU\_GF & As ABU, but the abundance variation is mimicked by a
          variation of $\log gf$ (Sect.~\ref{sec:loggf}). \\
ABU2   & Abundance from the wings of a strong line
         (Sect.~\ref{sec:abu2}). \\
MOLEC  & Abundance from a molecular feature
         (Sect.~\ref{sec:molec}). \\
ISOR   & Isotopic ratio (Sect.~\ref{sec:isor}). \\
ROTAT  & $v\sin i$, with the abundance free or fixed
         (Sect.~\ref{sec:rotat}). \\
\hline
\end{tabular}
\end{table}

%=====================================================================
\section{Synthetic spectra}
\label{sec:synth}
%=====================================================================

\subsection{Synthetic spectra computed on the fly}
\label{sec:synth_k2}

In its default configuration, \salvador{} computes the synthetic
spectra with the SYNTHE code \citep{Kurucz1993b, Sbordone2004,
Kurucz2005}, through the Python code K2 (Mucciarelli et al., in
prep.), which manages both ATLAS9 and SYNTHE. For each star, K2
computes a one-dimensional, plane-parallel, LTE ATLAS9 model
atmosphere \citep{Kurucz1993, Kurucz2005} with the exact $T_{\rm
eff}$, $\log g$, and [M/H] of the star, adopting the ODFs of the
KOALA database \citep{Mucciarelli2026} for the chosen [$\alpha$/Fe]
and interpolating them in [M/H] when needed. The KOALA ODFs are
computed with the solar mixture by \citet{Caffau2011} and
\citet{Lodders2010} and are available for $\xi=0$, 1, 2, 4, and
8\,km\,s$^{-1}$ \citep{Mucciarelli2026}. \todo{confirm with A.M.:
which ODF microturbulence is adopted by K2 for a given $\xi$} Each
model atmosphere is computed only once and reused for all the
features of the star and for all the following runs.

SYNTHE computes the emergent spectrum at a resolving power of
600\,000 in a spectral window centered on the feature, whose width is
chosen by the user (a few \AA\ are usually enough) and is
automatically enlarged when the fitting region requires it. The
atomic and molecular line lists are those of the Kurucz/Castelli
database\footnote{\url{https://wwwuser.oats.inaf.it/castelli/linelists.html}}
\todo{or other line lists: specify those adopted by default in K2}.
The line lists of TiO and H$_2$O can be included on request. The
unbroadened synthetic spectra are convolved by \salvador{} with a
Gaussian instrumental profile with the resolving power of the feature
(Sect.~\ref{sec:resolution}), using the PyAstronomy package
\citep{Czesla2019}, and resampled on the wavelength grid of the
observed spectrum.

The variations of the abundances (or of the isotopic ratios) required
during the analysis are passed to K2 together with the variations
given by the user in the file of abundance variations
(Sect.~\ref{sec:overview}). Each unbroadened synthetic spectrum is
stored with a name that encodes all the parameters that define it
(atmospheric parameters, chemical mixture, spectral window, and all
the abundance variations), so that a synthetic spectrum is computed
only once and reused whenever the same spectrum is needed, either in
the same run or in subsequent ones.

As discussed in Sect.~\ref{sec:intro}, the TiO and H$_2$O line lists
allow, in principle, the analysis of M giants, down to the lower
boundary of the KOALA grid, $T_{\rm eff}=3750$\,K
\citep{Mucciarelli2026}, at the cost of a dramatic increase of the
computing time of each synthetic spectrum.

\subsection{Grids of precomputed synthetic spectra}
\label{sec:synth_grid}

Alternatively, \salvador{} can work with a grid of precomputed,
continuum-normalized, unbroadened synthetic spectra, provided by the
user and calculated with any spectral synthesis code. For each feature,
\salvador{} selects the synthetic spectra with exactly the atmospheric
parameters, [M/H], and [$\alpha$/Fe] of the star, computed with
different abundances of the element under analysis, and covering the
wavelength of the feature. In this case, no new synthetic spectrum can
be computed during the analysis. Therefore, the abundances that can be
measured are limited to the range covered by the grid, and the features
whose $\chi^2$ minimum falls at the edge of the grid are rejected in the
automatic mode. The EWs are measured directly on the best-fit
synthetic spectrum, while $\chi$ and $\log gf$ of the transitions are
not available. The uncertainties due to the atmospheric parameters
(Sect.~\ref{sec:errors_params}) can be computed only if the grid also
includes the synthetic spectra with the perturbed parameters.

%=====================================================================
\section{Analysis of the spectral features}
\label{sec:method}
%=====================================================================

\subsection{Continuum normalization and local wavelength alignment}
\label{sec:norm}

The observed spectrum, read from a FITS or ASCII file, is cleaned of
the pixels with non-finite values and shifted to the rest frame with
the RV provided by the user. Three options are available for the
continuum normalization. The spectrum can be used as it is, if already
normalized; it can be rigidly rescaled to a synthetic spectrum
computed with the parameters of the star, by dividing it by the median
observed flux at the wavelengths where the synthetic spectrum is within
0.5\% of its maximum (i.e., at the continuum); or it can be normalized
by using a set of continuum windows provided by the user. In the latter
case, the pixels of each window that deviate by more than $5\sigma$
from the median (with $\sigma$ estimated from the median absolute
deviation) are removed, splitting the window into clean sub-windows,
the mean flux in each sub-window defines a continuum point, and the
continuum is obtained by interpolating these points with a cubic spline
(or linearly when fewer than four points are available). Since in
cool or metal-rich stars the continuum windows are rarely free of weak
absorption features, the normalized spectrum is then multiplied by the
mean level of the synthetic spectrum of the star in the same windows,
interpolated in the same way. This procedure takes into account the
pseudo-continuum predicted by the synthesis and avoids an
underestimate of the continuum level in crowded spectral regions.

Residual wavelength offsets between the observed and the synthetic
spectra, due for instance to uncertainties in the RV or in the
wavelength calibration, can be corrected locally. The center of the
observed feature is measured as the mean between the center of a
Gaussian fit and the midpoint of the two wavelengths at which the
profile crosses its half depth. If the offset between this center and
the laboratory wavelength is smaller than a tolerance set by the user
(in pixels), the observed spectrum is shifted accordingly for the
analysis of that feature.

\subsection{Spectral resolution}
\label{sec:resolution}

The resolving power adopted to broaden the synthetic spectra can be
specified for each feature in the line list. This allows the user to
account for the variation of the resolving power with wavelength, or
among different setups of the same instrument. Moreover, \salvador{}
can estimate the resolving power directly from each observed feature.
The full width at half maximum of the observed profile,
FWHM$_{\rm obs}$, is measured as the mean of the values obtained from a
Gaussian fit and from the half-depth crossings of a spline
interpolation of the profile. The same measurement is performed on the
unbroadened synthetic spectrum, after isolating the profile of the
main transition from the adjacent ones, providing the intrinsic width
FWHM$_{\rm syn}$, which includes the thermal, microturbulent, and
damping broadening. Assuming that the instrumental profile is Gaussian,
and approximating the intrinsic profile as Gaussian as well, the
instrumental width and the corresponding resolving power are
\begin{equation}
{\rm FWHM}_{\rm ins} = \sqrt{{\rm FWHM}_{\rm obs}^2 -
{\rm FWHM}_{\rm syn}^2}\,, \qquad
R = \frac{\lambda}{{\rm FWHM}_{\rm ins}}\,.
\label{eq:resolution}
\end{equation}
The estimated $R$ replaces the input value only if the relative
difference between the two is larger than 2\% and smaller than a
threshold chosen by the user, in order to protect the analysis
against spurious estimates (e.g., in weak, noisy, or blended
features). Since no macroturbulent broadening is included in the
synthesis, the resolving power estimated in this way also absorbs the
macroturbulence of the star, providing an effective Gaussian
broadening kernel for the feature. 
\subsection{Abundance determination}
\label{sec:abu}

For a given feature, the abundance of the element is derived by
minimizing
\begin{equation}
\chi^2(A) = \sum_{i}
\frac{\left[F_{{\rm obs},i} - F_{{\rm syn},i}(A)\right]^2}
     {F_{{\rm syn},i}(A)}\,,
\label{eq:chi2}
\end{equation}
where $F_{{\rm obs},i}$ and $F_{{\rm syn},i}(A)$ are the observed and
synthetic normalized fluxes at the $i$-th pixel of the fitting region,
and $A$ is the abundance of the element, $A({\rm X})=\log(N_{\rm
X}/N_{\rm H})+12$. For a spectrum dominated by photon noise, the
variance of the normalized flux at each pixel is $\sigma_i^2\simeq
F_i/({\rm S/N})_{\rm c}^2$, where $({\rm S/N})_{\rm c}$ is the
signal-to-noise ratio in the continuum. Therefore, Eq.~(\ref{eq:chi2})
corresponds, apart from the constant factor $({\rm S/N})_{\rm c}^2$,
which does not affect the position of the minimum, to the $\chi^2$ of a
photon-noise-limited spectrum, with the variance of each pixel
estimated from the model (Pearson's $\chi^2$). Additional noise sources,
such as read-out noise and sky subtraction, are not taken into
account. The limits of the fitting region are given by the
user or determined automatically from the observed profile.

The synthetic spectra needed to bracket the $\chi^2$ minimum are
computed adaptively, to keep their number as small as possible
(Fig.~\ref{fig:flowchart}). The first synthetic spectrum is computed with
the initial abundance given in the line list. The integrated line
depression, $\int(1-F)\,{\rm d}\lambda$, of the observed feature is
compared with that of the synthetic one, to establish whether the
abundance must be increased or decreased. A second synthetic spectrum
is then computed with an abundance variation of 0.4\,dex in that
direction. If the observed line depression lies between the two
synthetic ones, a third synthetic spectrum is computed at the
intermediate abundance; otherwise, two synthetic spectra are computed
at 0.6 and 0.8\,dex from the initial abundance. In both cases, three
synthetic spectra separated by 0.2\,dex are available. The $\chi^2$ of
each of them is computed and, if the minimum is not bracketed, the set
of three synthetic spectra is moved in steps of 0.2\,dex towards the
decreasing $\chi^2$, until the minimum is bracketed or the abundance
differs from the initial one by more than a threshold chosen by the user
(3\,dex by default), in which case the feature is rejected.

Once the minimum is bracketed, a fine grid of synthetic spectra with a
step of 0.01\,dex is built by linear interpolation in abundance of the
synthetic fluxes, pixel by pixel, between the synthetic spectra
adjacent to the one with the lowest $\chi^2$, and the abundance that
minimizes Eq.~(\ref{eq:chi2}) on this grid is adopted. The spacing of
0.2\,dex between the computed synthetic spectra is a compromise between
the number of syntheses and the accuracy of the linear interpolation of
the flux, which degrades for larger spacings, in particular for
saturated lines. Finally, the synthetic spectrum with the best-fit
abundance is computed directly with SYNTHE, together with two synthetic
spectra at $\pm0.2$\,dex, which are shown in the plot of the fit
(Fig.~\ref{fig:fit_example}).

After the abundance is determined, the EW of the feature is computed by
K2 from a synthetic spectrum that includes only the transition under
analysis (and its hyperfine or isotopic components) with the best-fit
abundance. Alternatively, all the transitions of the same species
falling in the feature can be included, as required for instance by the
Li resonance doublet. The EW is not used to derive the abundance, but
it provides the reduced EW, $\log({\rm EW}/\lambda)$, which is used,
together with $\chi$, to select the features entering the final
abundance (Sect.~\ref{sec:final_abu}) and to build the diagnostic plots
of Fe.

\begin{figure}
\centering
\fbox{\parbox[c][5cm][c]{0.9\hsize}{\centering
\todo{Fig. 2: example of fit of an Fe~I line (tutorial star
1500138 of M\,54, GIRAFFE HR18): observed spectrum, best fit,
$\pm0.2$\,dex synthetic spectra, continuum windows, and $\chi^2$ as a
function of $A$(Fe)}}}
\caption{Example of the fit of a Fe\,{\sc i} line.}
\label{fig:fit_example}
\end{figure}

\subsection{Abundances from the variation of the oscillator strength}
\label{sec:loggf}

When the abundance of an element is varied only in the spectral
synthesis, the thermal and pressure structure of the model atmosphere,
the electron pressure, and the continuous opacity remain those computed
for the original chemical composition. This approach does not introduce
significant errors as long as the abundance variation has a negligible
impact on the opacity and on the ionization structure of the atmosphere
\citep{Mucciarelli2026}. However, the spectral synthesis code
recomputes the ionization and molecular equilibria with the new
abundance, so that the synthetic spectrum is not computed with the
same chemical composition as the model atmosphere. This inconsistency
can become relevant for large abundance variations of the elements
that provide most of the free electrons, such as Mg, which is the main
electron donor in the atmospheres of FGK stars \citep{Mucciarelli2026},
or of the elements involved in the molecular equilibrium, such as C,
N, and O.

\salvador{} offers the possibility to avoid this inconsistency. In LTE,
the line absorption coefficient of a transition is proportional to the
product $gf\,N_{\rm X}$ of its oscillator strength and of the number
density of the absorbing element \citep[e.g.,][]{Gray2005}. Therefore,
as far as the line opacity is concerned, a variation $\Delta$ of
$\log gf$ of the transitions of the species under analysis is
equivalent to a variation $\Delta$ of $A$(X). With the keyword
ABU\_GF, the synthetic spectra are computed by varying the $\log gf$ of
these transitions instead of the abundance, while the chemical
composition adopted in the synthesis is the same as that of the model
atmosphere. The abundance is then $A({\rm X})=A_{\rm model}({\rm
X})+\Delta$, where $A_{\rm model}({\rm X})$ is the abundance of the
element in the model atmosphere and $\Delta$ is the best-fit variation
of $\log gf$. In this way, the synthetic spectrum is fully consistent
with the adopted model atmosphere. We note that neither approach is
equivalent to the computation of a new model atmosphere with the
derived abundance, which would require the opacity sampling technique
of ATLAS12 \citep{Kurucz2005} or new ODFs; when the derived abundance
differs significantly from that of the model atmosphere, the analysis
should be iterated with a more appropriate chemical mixture.

\subsection{Strong lines: abundances from the line wings}
\label{sec:abu2}

The core of strong, saturated lines is almost insensitive to the
abundance and forms in the outermost layers of the atmosphere, where
the assumptions of LTE and of a photospheric temperature structure are
less reliable, and where the ATLAS9 models may not reach full
convergence \citep{Mucciarelli2026}. On the other hand, the wings of
these lines, dominated by collisional damping, are formed in deeper
layers and remain sensitive to the abundance. With the keyword ABU2,
the $\chi^2$ of Eq.~(\ref{eq:chi2}) is computed over two separate regions
covering the blue and red wings of the line, excluding the core
(Fig.~\ref{fig:abu2}). The rest of the analysis is the same as described
in Sect.~\ref{sec:abu}. Since the damping wings depend on the gas
pressure, the abundances derived from them are sensitive to $\log g$
and to the chemical mixture of the model atmosphere
\citep{Mucciarelli2026}, and an accurate estimate of the uncertainty
due to $\log g$ is particularly relevant for these features
(Sect.~\ref{sec:errors_params}).

\begin{figure}
\centering
\fbox{\parbox[c][4cm][c]{0.9\hsize}{\centering
\todo{Fig. 3: the same strong line fitted with ABU (core) and ABU2
(wings), as in Fig. 8 of the manual}}}
\caption{Comparison between the fit of a strong line performed on
the core (keyword ABU, left) and on the wings (keyword ABU2, right).}
\label{fig:abu2}
\end{figure}

\subsection{Molecular features}
\label{sec:molec}

With the keyword MOLEC, the abundance is derived from a molecular
feature (e.g., CH, C$_2$, or CN bands) by varying the abundance of
the element under analysis, as for the atomic lines. The abundances of
the other elements forming the molecule can be fixed through the file
of abundance variations, for instance the C abundance when N is derived
from CN features. For these features, the EW, $\chi$, and $\log gf$ are
not computed.

\subsection{Isotopic ratios}
\label{sec:isor}

With the keyword ISOR, \salvador{} measures the ratio between the
abundances of two isotopes of the same element, such as
$^{12}$C/$^{13}$C from CN features or $^{6}$Li/$^{7}$Li from the Li
resonance doublet. The procedure is the same as for the abundances, but
the synthetic spectra are computed with different isotopic ratios,
with steps of 10 between the computed synthetic spectra (20 for the
first step), and the fine grid has a step of 1. The ratio is not allowed to go below unity. The
elemental abundance adopted in the synthesis is fixed, and should be
derived beforehand, for instance from features insensitive to the
isotopic ratio. The isotopic fractions of elements with more than two
isotopes can be set in the file of abundance variations.

\subsection{Projected rotational velocity}
\label{sec:rotat}

With the keyword ROTAT, \salvador{} derives the projected rotational
velocity, $v\sin i$, from a feature, either simultaneously with the
abundance of the element or at a fixed abundance. A grid of synthetic
spectra is computed with SYNTHE for abundances in steps of 0.1\,dex
within a range set by the user around the initial abundance (a single
abundance if it is kept fixed), and convolved with the instrumental
profile. Each of them is then broadened with the rotational profile of
\citet{Gray2005}, with a linear limb-darkening coefficient of 0.6, as
implemented in PyAstronomy \citep{Czesla2019}, for $v\sin i$ values
from a minimum to a maximum value (by default 0 and 60\,km\,s$^{-1}$)
in steps of 10\,km\,s$^{-1}$. Since the rotational broadening is
applied to the synthetic spectra already computed, a finer sampling in
$v\sin i$ does not require any additional synthesis. The $\chi^2$ of
Eq.~(\ref{eq:chi2}) is computed over this grid, and around its minimum
a fine grid with steps of 1\,km\,s$^{-1}$ (within
$\pm20$\,km\,s$^{-1}$) and 0.01\,dex (within $\pm0.2$\,dex) is built by
cubic interpolation (linear when fewer than four nodes are available
along an axis). The minimum of $\chi^2$ on the fine grid provides
$v\sin i$ and the abundance. Their uncertainties are estimated with a
Monte Carlo simulation, by adding to the best-fit spectrum Gaussian
noise with a dispersion equal to that of the residuals of the fit, and
repeating the search for the minimum (100 realizations by default).
As no macroturbulent broadening is included in the synthesis, the
derived $v\sin i$ includes any additional broadening not accounted for
by the instrumental profile; values comparable to the expected
macroturbulent velocity should therefore be regarded as upper limits.

\begin{figure}
\centering
\fbox{\parbox[c][4cm][c]{0.9\hsize}{\centering
\todo{Fig. 4: example of a ROTAT fit: best fit and spectra at
$\pm10$\,km\,s$^{-1}$, and $\chi^2$ contours in the
$A$(X)--$v\sin i$ plane with the Monte Carlo solutions}}}
\caption{Example of the simultaneous determination of $v\sin i$ and
abundance.}
\label{fig:rotat}
\end{figure}

\subsection{Interactive mode}
\label{sec:interactive}

In the interactive mode, the fit of each feature is displayed and the
user can accept it, reject it, or modify its ingredients: rescale the
normalized spectrum by a constant factor, shift it in wavelength,
change the limits of the fitting regions, change the resolving power,
add or remove continuum windows (the spectrum is then normalized
again), or move the set of synthetic spectra by one step towards the
decreasing $\chi^2$. After each change the fit is repeated, and the
initial configuration can be restored at any time. This mode is
especially useful for the analysis of a few stars or of critical
features (e.g., weak or blended lines, or lines in crowded spectral
regions), and to define the fitting regions to be adopted in the
automatic analysis of large samples.

\subsection{Final abundances}
\label{sec:final_abu}

For each star and species, the abundances of the individual features
are combined into a mean value. The features can be selected according
to their $\chi$ and reduced EW, within limits chosen by the user, to
exclude, for instance, strong lines that are more sensitive to $\xi$
and to the damping constants, or features too weak to be reliably
measured. A $\sigma$-clipping is then applied, adopting the median and
the median absolute deviation as estimators of the center and of the
dispersion, when at least three features are available. [Fe/H] is
derived from the Fe\,{\sc i} lines, and the abundance ratios of the
other species are provided as [X/Fe]. The absolute abundances $A$(X)
and the [X/H] ratios are also provided.

Different spectra of the same star, such as those obtained with
different setups of a multi-object spectrograph, can be analyzed
separately and their results merged afterwards, recomputing the final
abundances and their uncertainties with the same procedure.

%=====================================================================
\section{Uncertainties}
\label{sec:errors}
%=====================================================================

\subsection{Internal uncertainties}
\label{sec:errors_int}

For the species measured from fewer than ten features, the uncertainty
of each feature due to the photon noise is estimated through a Monte
Carlo simulation. The signal-to-noise ratio (S/N) is measured within
$\pm20$\,\AA\ from the feature, either in the continuum windows or,
when these are not available, on the flux after a $3\sigma$-clipping
that removes the absorption features. Gaussian noise with a dispersion
equal to (S/N)$^{-1}$ is added to the best-fit synthetic spectrum and
the abundance is derived again from the fine grid of
Sect.~\ref{sec:abu}. The procedure is repeated 200 times and, after
discarding the 5\% of the highest and lowest values, the standard
deviation of the derived abundances is adopted as the uncertainty. The
random numbers are generated from a seed computed from the name of the
spectrum, the wavelength, and the species of the feature, so that the
results are fully reproducible, independently of the order in which
the features are analyzed.

The internal uncertainty of the mean abundance of a species is computed
as $\sigma/\sqrt{N}$, where $\sigma$ is the standard deviation of the
abundances of the $N$ features that survived the selection. When only
one feature is available, the Monte Carlo uncertainty is adopted.

\subsection{Uncertainties due to the atmospheric parameters}
\label{sec:errors_params}

When the uncertainties of the atmospheric parameters are provided, the
analysis of each feature is repeated by varying, one at a time, $T_{\rm
eff}$, $\log g$, and $\xi$ by their uncertainties, while keeping the
other parameters fixed. New model atmospheres are computed by K2 for
the perturbed parameters. For each feature and parameter $p$, the
difference $\delta^{p}=A({\rm X})_{p+\Delta p}-A({\rm X})$ is stored
with its sign. The variation of $\log g$ is performed for all the
features. To save computing time, the variation of $\xi$ is performed
only for the features
with EW larger than 50\,m\AA, which are the most sensitive to this
parameter, and for those whose EW is not computed (molecular features,
line wings, and isotopic ratios).

For each species, the mean values of $\delta^{T_{\rm eff}}$,
$\delta^{\log g}$, and $\delta^{\xi}$ over the selected features are
computed. The total uncertainties of [Fe/H] and [X/Fe] are then
\begin{equation}
\sigma_{\rm [Fe/H]} = \sqrt{\frac{\sigma_{\rm Fe}^2}{N_{\rm Fe}} +
(\delta^{T_{\rm eff}}_{\rm Fe})^2 + (\delta^{\log g}_{\rm Fe})^2 +
(\delta^{\xi}_{\rm Fe})^2}\,,
\label{eq:err_feh}
\end{equation}
\begin{equation}
\begin{split}
\sigma_{\rm [X/Fe]} = \Bigg[\frac{\sigma_{\rm X}^2}{N_{\rm X}} &+
\frac{\sigma_{\rm Fe}^2}{N_{\rm Fe}} +
(\delta^{T_{\rm eff}}_{\rm X}-\delta^{T_{\rm eff}}_{\rm Fe})^2 \\
&+ (\delta^{\log g}_{\rm X}-\delta^{\log g}_{\rm Fe})^2 +
(\delta^{\xi}_{\rm X}-\delta^{\xi}_{\rm Fe})^2\Bigg]^{1/2},
\end{split}
\label{eq:err_xfe}
\end{equation}
where the internal terms $\sigma^2/N$ are replaced by the squared
Monte Carlo uncertainty when only one feature is available. Since the
differences are taken with their sign, the uncertainties due to the
atmospheric parameters partially cancel out in [X/Fe] for the species
that respond to the parameters as Fe does, while they add up for those
that respond in the opposite way, as some ionized species with respect
to $T_{\rm eff}$. The uncertainties of [X/H] are computed in the same
way, without the terms relative to Fe. This approach neglects the
covariances among the atmospheric parameters and should therefore be
regarded as an upper limit to the real uncertainty.

%=====================================================================
\section{Implementation and performance}
\label{sec:implementation}
%=====================================================================

\salvador{} is written in Python ($\geq3.10$) and relies on NumPy
\citep{Harris2020}, SciPy \citep{Virtanen2020}, pandas
\citep{McKinney2010}, Matplotlib \citep{Hunter2007}, Astropy
\citep{Astropy2013, Astropy2018, Astropy2022}, specutils
\citep{Specutils}, and PyAstronomy \citep{Czesla2019}. It is
distributed as a Python package under the MIT license, and it can be
installed with \texttt{pip} from the GitHub repository, together with
a tutorial and a user manual. The code is organized in modules
according to their task (input and output, normalization, synthesis,
fitting, rotation, uncertainties, and final abundances), and all the
settings of a run are read from a single input file, checked at the
start of the run.

Most of the computing time is spent in the calculation of the
synthetic spectra. To reduce it, \salvador{} (i) runs K2 within its own
Python process; (ii) passes to SYNTHE, for each feature, only the
portion of the line lists within 50\,\AA\ from the spectral window,
which provides synthetic spectra identical to those computed with the
complete line lists; (iii) reuses the model atmospheres and the
synthetic spectra already computed (Sect.~\ref{sec:synth_k2}); and
(iv) in the automatic mode, analyzes several features in parallel,
each in a separate process with its own working directory. The model
atmospheres needed by all the features are computed before the parallel
analysis starts, and the results are written in the same order as in
the line list, so that the output is identical to that of a sequential
run. As an example, the analysis of 25 lines in the GIRAFFE HR18
spectrum of a giant star of M\,54 ([Fe/H]$\,=-1.54$\,dex), with the
continuum normalized through continuum windows and the synthetic
spectra computed in windows of 5\,\AA, required 3.4\,min with a single
process and 52\,s with seven parallel processes (a speed-up of a
factor of 3.9), on a laptop with an Intel Core Ultra 7 258V processor
and 32\,GB of RAM. \todo{confirm: synth\_out initially empty (all
models and synthetic spectra computed from scratch)? errors due to the
atmospheric parameters computed?} The speed-up is lower than the number
of processes because the total time is set by the slowest features and
because the model atmospheres are computed before the parallel analysis
starts.

The reproducibility of the results is ensured by the seeded Monte Carlo
simulations (Sect.~\ref{sec:errors_int}) and by a regression tool that
compares the results of a run, value by value, with those of a
reference run. A second tool verifies that the faster calls to K2
provide synthetic spectra and EWs identical to those of the standard
calls. A suite of automatic tests, which replaces K2 with the synthetic
spectra of reference runs, is executed at every change of the code.

%=====================================================================
\section{A new version of GALA}
\label{sec:gala}
%=====================================================================
\todo{to be written with A. Mucciarelli}

%=====================================================================
\section{Validation}
\label{sec:validation}
%=====================================================================
\todo{synthetic spectra and benchmark stars}

%=====================================================================
\section{Conclusions}
\label{sec:conclusions}
%=====================================================================

In this paper we presented \salvador{}, a Python code for the
determination of chemical abundances through spectral synthesis,
designed for the line-by-line analysis of FGK stars in the optical and
infrared domains, from the low resolution of integral-field
spectrographs such as MUSE ($R\sim3000$) to the very high resolution
of ESPRESSO. For each feature, \salvador{} minimizes the $\chi^2$
between the observed spectrum and a set of synthetic spectra that
differ only in the quantity to be measured. The synthetic spectra are
computed on the fly with ATLAS9 and SYNTHE through K2, adopting the
KOALA opacity distribution functions \citep{Mucciarelli2026}, or taken
from a grid of precomputed spectra calculated with any spectral
synthesis code. An adaptive search keeps the number of syntheses
small, and the minimum is refined on a grid with a step of 0.01\,dex
obtained by interpolation of the synthetic fluxes.

Besides the standard abundance analysis, \salvador{} combines within a
single framework several features that are rarely available together:
\begin{itemize}
\item the abundance variation can be mimicked by a variation of the
$\log gf$ of the transitions, keeping the chemical composition of the
synthesis consistent with that of the model atmosphere;
\item the resolving power can be estimated from each observed feature,
accounting for the intrinsic width of the line;
\item the abundances of saturated lines can be derived from their
damping wings only, excluding the line core;
\item isotopic ratios and the projected rotational velocity (with the
abundance either free or fixed) can be measured;
\item the uncertainties due to the photon noise are estimated through
reproducible Monte Carlo simulations, and those due to the atmospheric
parameters are computed line by line and propagated with their sign
into the abundance ratios.
\end{itemize}
The interactive mode allows full control of every fit, while the
automatic mode, with the analysis of several features in parallel,
makes the code suitable for large spectroscopic samples.

\todo{Summary of the validation on synthetic spectra and benchmark
stars (Sect.~\ref{sec:validation}): accuracy and precision of the
abundances, of $v\sin i$, and of the isotopic ratios.}

\salvador{} does not derive the atmospheric parameters, which must be
provided by the user. The new version of GALA presented in
Sect.~\ref{sec:gala} fills this gap, and the two codes together provide
a complete chain for the chemical analysis of stellar spectra, from
the atmospheric parameters to the final abundances and their
uncertainties. Future developments of \salvador{} will include the
convolution of the synthetic spectra with non-Gaussian instrumental
profiles and a continuum normalization based on a polynomial fit of
the ratio between the observed and the synthetic spectra.
\salvador{} is publicly available on GitHub, together with a user
manual and a tutorial.


\begin{acknowledgements}
BLABLA
\end{acknowledgements}

\bibliographystyle{aa}
\bibliography{salvador}

\end{document}
